\chapter{Introduction}

Few theories have reshaped modern technologies and thus our day to day life as profoundly as quantum mechanics when it was discovered in the early 1900s. The examination of the black-body radiation of celestial bodies such as our sun and later formulation of the Planck's law by Max Planck, explaining the quantized energy exchange between a black-body and a radiation field, laid the ground work for many theories and technologies. Among these is the laser laser which is based on Einstein's theory on stimulated emission \cite{EinsteinLaser}, both used in day to day life (e.g. the barcode scanner at the register, or scanner within a computer mouse) as well as research (e.g. Sisyphus cooling \cite{SysiphusCoolingWineland} or trapping of atoms in a MOT \cite{MOTPaper}). Another technology of great importance is the atomic clock, which is based on Norman Ramsey's work \cite{RamseyNobelPrize}, allowing for the clear definition of a single second \cite{RamseySecondDefinition}. One piece of technology, which for the average person, stands above all of those, is that of the transistor, specifically the MOSFET transistor. The MOSFET transistor is the most used type of transistor found in modern electronics \cite{MOSFET} that either allows or denies the passage of electrical current. By either blocking the current or letting it pass through the component, one forms the foundation for all modern electronics, namely that of a \textit{bit}. Interestingly enough, the theory which explains the passage of current through such a device, that being the electronic band structure model, is again rooted in quantum mechanics.\cite{Hunklinger} \\

In an attempt to predict the following years of chip evolution, Gordon Moore made the bold assumption in 1965 that with each coming year, the number of components crammed on an integrated chip will double \cite{MooresLaw}. This statement would later on be known as \textit{Moore's Law}, and would pave the way for chip manufacturers to keep on increasing the amount of transistors in their chips, in order to subconsciously keep up with the predictions made by Moore. Back in the day, with the way technology was evolving, this enthusiasm towards simply increasing the amount of transistors by following a "plenty of room at the bottom"-mentality from the likes of Richard Feynman \cite{feynman1992plenty} and thus achieving an increase in computing performance, would eventually hit fundamental limits. In order to pack more transistors onto the same sized computer chip, the transistors themselves would need to be shrunk down. However in recent years, this approach of increasing computing performance has led to the realization that simply going smaller, does not yield the hoped-for performance increase. \cite{MooresLawBad} \\

As the tasks one concerns themselves with become much more intricate, like finding the most efficient delivery route for a logistics company\cite{TravellingSalesmanProblem}, or accurately simulating small subsystems \cite{FeynmanSimulatingPhysicsWithQuantumComputers,QCMedicine,BlochSimulation}, it's clear that classical computation has hit a limit and a new kind of computation platform is needed. As Feynman points out in his talk "Simulating Physics with Computers":
\begin{quote}
	``Nature isn't classical, dammit, and if you want to make a
	simulation of nature, you'd better make it quantum mechanical,
	and by golly it's a wonderful problem, because it doesn't look
	so easy.''\cite{FeynmanSimulatingPhysicsWithQuantumComputers}
\end{quote}
This has paved the way for a new form of computing, called \textit{quantum computing}, which instead of using fixed bits that can either be a 1 or a 0, this new type of computer uses the so called \textit{qubits}, which in itself is a two-state quantum-mechanical system, and can be placed in any \textit{superposition} state $\ket{\psi}=\alpha\ket{0} + \beta\ket{1}$. That is, until a measurement is performed and the state collapses back to being either $\ket{0}$ or $\ket{1}$ again. What makes the qubit interesting however, is the relationship multiple qubits can have with one another, and that they can be \textit{entangled} with one another, thus not allowing for description of one qubit state, without the state of another qubit. By leveraging these two fundamental quantum mechanical principles, many protocols have been implemented that in theory should outperform their classical counterparts. Be it in the realm of encryption where certain quantum algorithms now promise an exponential speedup \cite{ShorrsAlgorithm}, in quantum key distribution where the laws of quantum mechanics enable fundamentally secure communication \cite{BB84Protocol}, or in the transfer of quantum states between distant nodes via quantum teleportation \cite{QuantumTeleportation}.\\

However, as time went on, and different hardware platforms were investigated that could serve as a foundation for a quantum computer, one problem (along many others), became apparent: scalability. As one tries to increase the amount of qubits used in a quantum computer, being able to keep external noise and errors that each component produces at bay, becomes more and more of a challenge\cite{ScalabilityProblemQC}. In order to be able to scale these new systems up, it was proposed that instead of building one large quantum computer, where one would try to keep errors at bay via error correction algorithms integrated into the fundamental structure of the computer \cite{GoogleErrorCorrection}, one would instead split the computation horsepower up into multiple smaller decentralized pieces.\cite{DuanKimble2004,QuantumInternetKimble2008} Each of these pieces would then represent a single quantum node, that is able to create entanglement between a stationary atom and a mediating qubit, transfer the state of the mediating qubit onto the stationary one, and also process information through the use of quantum logic gates. By linking many of these quantum nodes one aims to overcome size-scaling and error-correction problems that would limit the size of machines for quantum processing.\cite{QuantumInternetKimble2008,ScalableQCArchitecture_silicon_based}. This network of nodes would then appropriately be named a \textit{Quantum Internet}. One way of realizing such a quantum internet is via the strong coupling interaction of single photons and atoms in the setting of cavity quantum electrodynamics (\textit{cQED}). Here, each quantum node would store a set number of stationary atoms in between two highly reflective mirrors, which are able to be individually controlled, and interlinked with one another through another set of quantum channels, realized by photons that travel in between the nodes through optical fibers.\\

One of these cavity systems is the \textit{CavityX} systems at the Max Planck Institute for Quantum Optics which features a $^{87}Rb$ atom trapped in between two crossed fiber optic cavities, hence the letter X in the name. In the past this setup has been used in order to achieve atom-photon entanglement with a fidelity of 87\% \cite{HeraldedEntanglement}, successfully transfer the information from a photonic qubit onto an atomic qubit with a fidelity of of 94.7\%\cite{QuantumNetworkManuel}, and recently realize atom-atom entanglement between two nodes connected via a fiber going through the Munich city center all the way out to research campus in Garching while still achieving a fidelity of \notetome{include number and citation here}. However as mentioned before, in order to realize a full quantum internet, the last milestone is the realization of the processing nature of the quantum node part mentioned in the definition of the quantum internet by Jeff Kimble \cite{QuantumInternetKimble2008}. This is where this thesis comes into play. The goal of this thesis is to realize an atom-photon reflection gate using one of the two fiber cavities in the CavityX setup, that one being birefringent, while preparing the atom in either of the two clock states of the $\text{D}_2\text{-line}$ ground states, thus minimizing the influence of external factors like magnetic field fluctuations. Using a setup inspired by \cite{DuanKimble2004} and storing the information in the polarization of a photonic qubit, has the advantage of allowing for controlled gate operations between a flying photon and a resonantly coupled atom trapped inside of the cavity. With the realization of a quantum gate between a stationary atom and a flying photonic qubit, one would be able to reuse this setup to build further quantum nodes as such and thus completing the final requirement for a fully functional quantum network node, ready for integration into a Quantum Internet.\\

\notetome{Include the structure of thesis once Pau is done giving me his verdict}
